\documentclass{article}
\usepackage{import}
\usepackage{tcolorbox}
\usepackage{amsmath}
\usepackage{amssymb}
\usepackage{amsthm}
\usepackage{geometry}
\usepackage{enumitem}
\usepackage{pdfpages}
\usepackage{transparent}
\usepackage{booktabs}
\usepackage{bm}
\usepackage{tabularx}
\usepackage{multirow}
% \usepackage{minted}
\usepackage{xcolor}
\usepackage{setspace}
\usepackage[hidelinks]{hyperref}
\usepackage{pgfplots}
% \usepackage{soulpos}

\tcbuselibrary{theorems}
\tcbuselibrary{skins}

\setcounter{tocdepth}{6}
\setcounter{secnumdepth}{6}

\newcommand{\param}{\boldsymbol{\theta}}
\newcommand{\bs}{\boldsymbol}
\newcommand{\data}{\mathcal D}

% pgfplots settings
\pgfplotsset{compat=1.18}
\usepgfplotslibrary{external}

% tabularx Settings
\renewcommand\tabularxcolumn[1]{>{\centering\arraybackslash}m{#1}}

% Page Settings
\newgeometry{vmargin={15mm}, hmargin={15mm, 15mm}}
\setlength{\parindent}{0pt}
\setlength{\parskip}{0.7\baselineskip}%{0.8em plus 1pt minus 1pt}

\onehalfspacing

% hyperref Settings
\hypersetup{
    colorlinks=true,
    linkcolor=blue,
    filecolor=magenta,      
    urlcolor=cyan,
}

% inkscape Settings
\newcommand{\incfig}[1]{%
    \def\svgwidth{\columnwidth}
    \import{./figure/}{#1.pdf_tex}
}

% inline code block settings
\newtcbox{\mybox}[1][]{
    on line,
    arc=1pt, outer arc=2pt,
    colback=lightgray!50!white, colframe=lightgray!50!white,
    boxsep=0pt, left=0pt, right=-1pt, top=2pt, bottom=0pt,
    boxrule=0pt, #1
}

\NewDocumentCommand{\code}{v}{%
  \begingroup\ttfamily
  \codeulinner{#1}%
  \endgroup%
}

% Definition Box Styles
\newtcbtheorem[auto counter, number within=subsection]{defn}{Definition}{
% basic settings
    theorem name and number,
    separator sign={},
    description delimiters={(}{)},
    enhanced jigsaw,
    sharp corners,
% box margins and rules
    boxrule=1pt,
    left=0.2cm,right=0.2cm,top=0.2cm,
    toptitle=0.1cm,
    bottomtitle=0.08cm,
% box colors
    colframe=white!25!black,
    colback=white,
    coltitle=white,
% box fonts
    fonttitle=\bfseries,fontupper=\normalsize
}{defn}

% Theorem Box Styles
\newtcbtheorem[auto counter, number within=subsection]{thm}{Theorem}{
% basic settings
    theorem name and number,
    separator sign={},
    description delimiters={(}{)},
    enhanced jigsaw,
    sharp corners,
% box margins and rules
    boxrule=1pt,
    left=0.2cm,right=0.2cm,top=0.2cm,
    toptitle=0.1cm,
    bottomtitle=0.08cm,
% box colors
    colframe=white!25!black,
    colback=white,
    coltitle=white,
% box fonts
    fonttitle=\bfseries,fontupper=\normalsize
}{thm}

% Lemma Box Styles
\newtcbtheorem[number within=tcb@cnt@thm]{lem}{Lemma}{
% basic settings
    theorem name and number,
    separator sign={},
    description delimiters={(}{)},
    enhanced jigsaw,
    sharp corners,
% box margins and rules
    boxrule=1pt,
    left=0.2cm,right=0.2cm,top=0.2cm,
    toptitle=0.1cm,
    bottomtitle=0.08cm,
% box colors
    colframe=white!50!black,
    colback=white,
    coltitle=white,
% box fonts
    fonttitle=\bfseries,fontupper=\normalsize
}{lem}

% Corollary Box Styles
\newtcbtheorem[number within=tcb@cnt@thm]{cor}{Corollary}{
% basic settings    
    theorem name and number,
    separator sign={},
    description delimiters={(}{)},
    enhanced jigsaw,
    sharp corners,
% box margins and rules
    boxrule=1pt,
    left=0.2cm,right=0.2cm,top=0.2cm,
    toptitle=0.1cm,
    bottomtitle=0.08cm,
% box colors
    colframe=white!50!black,
    colback=white,
    coltitle=white,
% box fonts
    fonttitle=\bfseries,fontupper=\normalsize
}{cor}

% Remark Box Styles
\newtcbtheorem[number within=subsection]{rem}{Remark}{
% basic settings
    theorem name and number,
    separator sign={},
    description delimiters={(}{)},
    enhanced jigsaw,
    sharp corners,
% box colors
    colframe=white!60!black,
    colback=white,
    coltitle=black,
% box margins and rules
    boxrule=1pt,
    left=0.2cm,right=0.2cm,top=0.2cm,
    toptitle=0.1cm,
    bottomtitle=0.08cm,
% box fonts
    fonttitle=\bfseries,fontupper=\normalsize
}{rem}

% Example Box Styles
% Remark Box Styles
\newtcbtheorem[number within=subsection]{exm}{Example}{
% basic settings
    theorem name and number,
    separator sign={},
    description delimiters={(}{)},
    enhanced jigsaw,
    sharp corners,
% box colors
    colframe=white!60!black,
    colback=white,
    coltitle=black,
% box margins and rules
    boxrule=1pt,
    left=0.2cm,right=0.2cm,top=0.2cm,
    toptitle=0.1cm,
    bottomtitle=0.08cm,
% box fonts
    fonttitle=\bfseries,fontupper=\normalsize
}{exm}


% Proof Box Styles

\title{Parameter Initialization}
\author{Hyoung Joon, Kim}
\date{\today}

\begin{document}
\maketitle
\tableofcontents
\newpage

\section{Introduction}
    We introduce some ways of initializing parameters. One may think like "Is it that important? Can't we just initialize it randomly?".
    But if we just randomly initialize it, we have two problems:
    \begin{enumerate}
        \item \textbf{Vanishing Gradient};
        \item \textbf{Exploding Gradient}.
    \end{enumerate}
    Actually this problems arise due to activation functions, not randomly intialized parameters.
    We explain it more on next section.
\section{Main Ideas}
    \subsection{Vanishing and Exploding Gradients}
        Let's look at the backpropagation equation (the chain rule).
        \[
            \frac{\partial f}{\partial x_1} =
            \frac{\partial f}{\partial x_n}\frac{\partial x_n}{\partial x_{n-1}}\dots\frac{\partial x_2}{\partial x_1}
        \]
        The values are multiplied on and on. The point is that if we use activation functions such as $\sigma(\dot)$, the sigmoid,
        the gradient values between 0 and 1 are keep multiplied. This will result in very small gradients if the network if very deep,
        and float arithmetic won't be able capture the value. This is the \textbf{vanishing gradient problem}.
        Even if the gradients being multiplied is fairly big enough, this may \textbf{lead to exploding gradient problem}.
        The forward pass also can suffer from very big or small values, since the values are kept multiplied and added by big or small values.

    \subsection{Xavier Intialization}
        \textbf{Xavier initialization}, or \textbf{Glorot initialization} (Glorot \& Bengio, 2010), is a way to initialize weights,
        which relaxes the problems shown before. This method is derived if we try to maintain the variances of gridients during backward pass
        and input vectors during forward pass constant. For more information, see Glorot \& Bengio, 2010.
        \begin{defn}{Xavier Initialization}{xavierInit}
            Let $n_{in}$ be the number of input dimension of the layer and let $n_{out}$ be the number of output dimension of the layer.
            The \textbf{Xavier normal initialization} is to pick the weight from following normal distribution:
            \[
            \mathcal{N}(0, \sigma^2),\ \text{where}\ \sigma^2 = \sqrt{\frac{2}{n_{in}+n_{out}}}.
            \]
            \textbf{Xavier uniform initialization} is to pick the weight from following uniform distribution:
            \[
            \mathcal{U}\left(-\sqrt{\frac{6}{n_{in}+n_{out}}},\ \sqrt{\frac{6}{n_{in}+n_{out}}}\right).
            \]
        \end{defn}
        The uniform distribution is the one proposed on the paper. The Xavier initialization works well on tanh and sigmoid function, but not on
        ReLU. The other method, \textbf{He initialization} (He et al., 2015), is told to work better on ReLU activations.

    \subsection{He Initialization}
        \begin{defn}{He Initialization}{heInit}
            Let $n_{in}$ be the input dimension of the layer. Than, we pick the weights from following normal distribution:
            \[
            \mathcal{N}\left(0,\ \frac{2}{n_{in}}\right).
            \]
        \end{defn}
        This method, like Xavier initialization, can be derived from trying to maintain the variances constant.
    
    \subsection{Put more methods if you learned other one!}

\end{document}